\noindent\textbf{Fleet Size.}
To evaluate the sensitivity of DVNDA to fleet size, we fixed the customer
scale at $n=20$ and varied the number of homogeneous vehicles as
$m\in\{4,5,6,7\}$. The $m=4$ results were taken from the original experiment.
For $m=5,6,7$, a separate DVNDA model was trained using the same architecture
and training configuration, while all other problem settings were unchanged.
Each model was evaluated on 100 test instances at four dynamic rates.
Table~\ref{tab:fleet_size} reports route Cost and QoS.

\begin{table*}[htbp]
\centering
\caption{Sensitivity of DVNDA to different fleet sizes on 20-customer instances. The $m=4$ row is taken from the original experiment; the other rows are newly trained and evaluated on 100 instances.}
\label{tab:fleet_size}
\small
\setlength{\tabcolsep}{3.5pt}
\begin{tabular}{c|cc|cc|cc|cc}
\toprule
\multirow{2}{*}{$m$} & \multicolumn{2}{c|}{10\%} & \multicolumn{2}{c|}{25\%} & \multicolumn{2}{c|}{50\%} & \multicolumn{2}{c}{75\%} \\
& Cost & QoS & Cost & QoS & Cost & QoS & Cost & QoS \\
\midrule
4 & 8.31 $\pm$ 1.22 & 100.00\% & 8.95 $\pm$ 1.30 & 100.00\% & 10.47 $\pm$ 1.53 & 100.00\% & 11.78 $\pm$ 1.44 & 100.00\% \\
5 & 9.33 $\pm$ 1.26 & 100.00\% & 9.76 $\pm$ 1.12 & 100.00\% & 11.04 $\pm$ 1.25 & 100.00\% & 11.87 $\pm$ 1.29 & 100.00\% \\
6 & 8.59 $\pm$ 1.30 & 100.00\% & 9.28 $\pm$ 1.20 & 100.00\% & 10.32 $\pm$ 1.19 & 100.00\% & 11.40 $\pm$ 1.28 & 100.00\% \\
7 & 8.75 $\pm$ 1.17 & 100.00\% & 9.49 $\pm$ 1.18 & 100.00\% & 10.46 $\pm$ 1.14 & 100.00\% & 11.37 $\pm$ 1.16 & 100.00\% \\
\bottomrule
\end{tabular}
\end{table*}


As the dynamic rate increased, Cost increased consistently for every fleet
size, while QoS remained at 100\% in all settings. The variation across fleet
sizes was moderate but non-monotonic: $m=6$ achieved the lowest Cost at the
10\%, 25\%, and 50\% dynamic rates, whereas $m=7$ was slightly better at 75\%.
The maximum Cost difference among the four fleet sizes decreased from 1.02 at
the 10\% dynamic rate to 0.51 at 75\%. These results indicate that DVNDA
maintains stable service quality and comparable routing performance under
moderate changes in fleet size.
